% appendix_experiments.tex — Appendix A-B for ask-factors NeurIPS 2026 paper
% Additional experimental details and results

\section{Implementation Details}\label{app:implementation}

\paragraph{Hyperparameters.}
Table~\ref{tab:hyperparams} provides the full hyperparameter configuration used across all AFM experiments.

\begin{table}[h]
\centering
\caption{\textbf{Hyperparameter settings for all AFM experiments.} Unless otherwise noted, all values are shared across the four AFM variants.}
\label{tab:hyperparams}
\small
\begin{tabular}{ll}
\toprule
\textbf{Hyperparameter} & \textbf{Value} \\
\midrule
Number of factors ($F$) & 5 (2 long-only, 3 long-short) \\
Lookback window ($L$) & 252 trading days \\
State dimension ($S$) & 16 \\
Hidden dimension (MLP) & 128 \\
Hidden dimensions (Residual MLP) & 128 $\to$ 64 \\
Activation function & GELU \\
Normalization & Layer normalization \\
Dropout rate & 0.3--0.5 \\
Optimizer & Adam \\
Learning rate & 0.01 \\
Batch size & 128 (time indices) \\
Training epochs & 150 (with early stopping) \\
Orthogonality penalty ($\lambda_{\text{orth}}$) & $10^{-5}$ to $10^{-2}$ \\
Loading regularization ($\lambda_\beta$) & L2 penalty \\
Weight regularization ($\lambda_w$) & L2 penalty \\
\midrule
\multicolumn{2}{l}{\textbf{Portfolio weight construction}} \\
Long-only factors & Softmax over scores \\
Long-short factors & Softmax-minus-softmax with gross leverage normalization \\
\midrule
\multicolumn{2}{l}{\textbf{Evaluation}} \\
Train/test split & Last 10\% of dates held out \\
Metric & Per-fund OOS $R^2$, averaged across funds \\
\bottomrule
\end{tabular}
\end{table}

\paragraph{Parameter counts.}
Table~\ref{tab:param_counts} reports the parameter counts for each model configuration. The parameter count varies primarily with the backbone architecture and the number of assets in the training universe (which determines the dimensionality of the ReturnEncoder input).

\begin{table}[h]
\centering
\caption{\textbf{Parameter counts by model configuration.}}
\label{tab:param_counts}
\begin{tabular}{lccr}
\toprule
\textbf{Model} & \textbf{Backbone} & \textbf{Universe} & \textbf{Parameters} \\
\midrule
AFM MLP All & MLP (128) & All 10 classes & 34.8M \\
AFM ResMLP All & ResMLP (128$\to$64) & All 10 classes & 69.8M \\
AFM MLP Large Cap & MLP (128) & Large cap only & 13.8M \\
AFM ResMLP Large Cap & ResMLP (128$\to$64) & Large cap only & 27.7M \\
Fama-French 4-Factor & N/A (OLS) & Large cap & 0 \\
\bottomrule
\end{tabular}
\end{table}

\paragraph{Software and hardware.}
% TODO: INSERT ACTUAL SOFTWARE VERSIONS AND GPU MODEL HERE
All AFM models are implemented in PyTorch. Training is performed on a single NVIDIA GPU. Each full training run (150 epochs) completes within several hours. The Fama-French baseline is implemented via ordinary least squares regression using standard linear algebra routines.

\paragraph{Characteristic preprocessing.}
Fund-level characteristics are extracted from SEC N-PORT filings and include two groups: (i) economic sector holdings (percentage of fund AUM allocated to each GICS sector) and (ii) asset class allocations (percentage in equities, fixed income, cash, alternatives, etc.). All characteristics are cross-sectionally normalized per date to have zero mean and unit variance across the fund universe.

\section{Additional Per-Category Results}\label{app:additional_results}

\paragraph{Median OOS $R^2$.}
Table~\ref{tab:median_r2} reports the median (rather than mean) OOS $R^2$ across funds. The median is more robust to outlier funds and provides a complementary view of typical model performance. The ranking of models is preserved: AFM MLP All achieves the highest median $R^2$ at 88.1\%, compared to 75.4\% for Fama-French.

\begin{table}[h]
\centering
\caption{\textbf{Median OOS $R^2$ (\%) across models.} The median is computed across all funds in each model's evaluation universe. Higher is better.}
\label{tab:median_r2}
\begin{tabular}{lc}
\toprule
\textbf{Model} & \textbf{Median OOS $R^2$} \\
\midrule
AFM MLP All & \textbf{88.1\%} \\
AFM ResMLP All & 85.4\% \\
AFM MLP Large Cap & 86.0\% \\
AFM ResMLP Large Cap & 81.5\% \\
Fama-French 4-Factor & 75.4\% \\
\bottomrule
\end{tabular}
\end{table}

\paragraph{Best epoch by configuration.}
Table~\ref{tab:best_epoch} reports the epoch at which each AFM variant achieves its best validation loss. All models converge well before the 150-epoch training budget, and the best epochs cluster in the 84--114 range.

\begin{table}[h]
\centering
\caption{\textbf{Best validation epoch by model configuration.}}
\label{tab:best_epoch}
\begin{tabular}{lc}
\toprule
\textbf{Model} & \textbf{Best Epoch} \\
\midrule
AFM MLP All & 84 \\
AFM ResMLP All & 100 \\
AFM MLP Large Cap & 112 \\
AFM ResMLP Large Cap & 114 \\
\bottomrule
\end{tabular}
\end{table}

% TODO: Add per-fund R² distribution histograms if available
% TODO: Add factor portfolio weight analysis (concentration, turnover) if available
% TODO: Add sensitivity to number of factors (3, 5, 7) if additional runs are conducted
